% Appendix F -- case-study details (moved from Section VII for length, 2026-10-12; text
% unchanged apart from joins). Every number is a macro of numbers.yaml filled by
% tools/gen_results.py from paper/sweeps/results and sremu validation/fm (or a definition
% macro); figure F12 comes from figures/src/f12_*.py. E-D is "model-based dimensioning, not
% validated". E-E: sremu examples/gnuradio_zmq, experiments/ee_gnuradio.py; results in
% sweeps/results/E-E (laptop WSL2 run of a working tree, to be repeated at the pinned commit)
% and E-E-laptop.
\section{Case-Study Details}\label{app:cases}

This appendix extends Section~\ref{sec:cases}.

\subsection{E-B: setup, cases and lab-like geometry}
\claim{C-EB-WHY}Three lessons stand out in E-B: an equalizer trained before a switch inside a
packet sees an error as large as the signal after it; a long linear transition spreads the
damage but does not remove it; and with a strong line of sight QPSK tolerates the wrong
configuration, so the timing of a command matters where the RIS path dominates or the margin
is small.

\claim{C-EB-SETUP}In the arc geometry each uncertain distance and angle takes the middle of its
known range, the two receivers stand at different angles at positions that are an engineering
choice (ENG), and $E_s/N_0$ at RX1 is set against the physical noise. A lab-like SIM geometry with one panel of three tiles is a supporting
case. Three assumptions apply to both geometries. The two states of the laboratory cells have
unit amplitude ($|\Gamma|=1$, ASSUMED), because their magnitudes were not measured. During a
transition the response follows the complex-linear path between the states, an engineering
choice (ENG). When the line of sight is blocked, the occluder removes the direct path and leaves
the RIS hops unattenuated; on the arc a real plate of that size would also shadow some RIS
rays, so this case is an idealization. The packets run at \numEbSps{} samples per symbol and
\numEbRateMsps~MS/s.

\claim{C-EB-TIMING}The slice of E-B2 shown uses guard intervals of \numEbGuardMs~ms, and on the
arc each controller message carries one command per tile. The arc run fixed three cases before
any point ran: the line of sight present at $E_s/N_0=\numEbEsnoDb$~dB, the same at
\numEbArcEsnoLowDb~dB as a sensitivity, and the line of sight blocked. With the line of sight
present at \numEbEsnoDb~dB QPSK tolerates the wrong configuration, although the EVM of the
affected packets rises. With a lead of \numEbLeadMinMs~ms, timed commands fail already at
$J=\numEbJitterOneMs$~ms (\numEbArcTimedHalfPerOneLosPct~\% with the line of sight present at
\numEbArcEsnoLowDb~dB and \numEbArcTimedHalfPerOneBlkPct~\% with it blocked); with
\numEbLeadMaxMs~ms they lose no packet at any jitter. \claim{C-EB-SERVED}On the served path of
Section~\ref{sec:wc-loop} an asap command became visible after a median of
\numBfourVisPfiftyMinMs--\numBfourVisPfiftyMaxMs~ms, with p99
\numBfourVisPninetynineMinMs--\numBfourVisPninetynineMaxMs~ms and rare outliers of up to
\numBfourVisMaxMs~ms; a timed command there is on time when its lead exceeds the command
latency plus the error of the client's estimate of virtual time.
\claim{C-EB-TICK}The clock of the tick-based semantics starts at $t=0$
and keeps one pending command per panel. Its losses hold in both cases of
Fig.~\ref{fig:F9}(b): with a timer period of the order of a packet, alignment with the frame
decides.

\claim{C-EB-LABLIKE}In the lab-like geometry (Fig.~\ref{fig:F9}(c), (d)), with the line of
sight blocked, the switch changes the channel of RX1 by
\numEbBlockedChangeDb~dB and \numEbBlockedRotDeg$^\circ$, and \numEbOneErrors{} payload bits
fail for the same switch time. Timed commands with the \numEbLeadMidMs~ms lead lose no packet
up to $J=\numEbJitterOneMs$~ms and \numEbTimedTwoPerThreePct~\% at $J=\numEbJitterMaxMs$~ms.
Asap commands lose \numEbAsapPerOnePct~\% at $J=\numEbJitterOneMs$~ms. Tick-based semantics
lose \numEbTickPerZeroPct~\% without jitter and up to \numEbTickPerMaxPct~\% with it. With the
line of sight present, no semantics loses a packet at \numEbEsnoDb~dB; this geometry was not
run at \numEbArcEsnoLowDb~dB.

\subsection{E-A: protocol record and time scale}
\claim{C-EA-SETUP}The access point sits \numEaApOffsetM~m off the user line, which keeps
the line-of-sight power finite everywhere on the line, and the partition loss gives the RIS
room to help. The
hypotheses and their ranges were in the emulator's design document before the sweep: GitHub's
push log (Section~\ref{sec:vv-lab}) records it in the initial push of the emulator repository.
The plan that adds the decision rules and the evaluation rule (SHA-256 \txtEaPlanSha\ldots{}
with LF line ends) was written before the sweep by the authors' record, but committed to the
paper repository after it (commit \txtEaPlanCommit), so no third party dates those parts. One
amendment, made before the full sweep after a short IQ run showed that evaluating each time
step at its start overestimates the power, moved the evaluation to the step's midpoint; the
plan discloses it and a short pipeline test seen before it.

\claim{C-EA-RESULT}For isotropic scattering the field and envelope correlations fall to one
half after \numClarkeFieldHalf{} and \numClarkeEnvHalf{} $\lambda/v$~\cite{Clarke1968}. The
rule-of-thumb coherence time $\sqrt{9/(16\pi)}/f_D=\numCoherenceRuleLambdaV\,\lambda/v$ lies
close to the E-A constant, but it is the geometric mean of $1/f_D$ and the envelope
half-correlation time, not a correlation level. The ratio of the two half-lives, \numEaRatio{}
(\numEaRatioLo--\numEaRatioHi), lies \numEaRatioOffsetPct~\% above the wavelength ratio
\numEaRatioPred; the corridor is the same in metres at both carriers, so it is not the same
geometry in wavelengths and exact scaling is not expected. The ranges fixed before the sweep
were \numEaHaOneLoS--\numEaHaOneHiS~s for H-A1 and \numEaHaTwoLo--\numEaHaTwoHi{} for H-A2, so
their support is a weak test. What E-A adds to the classical scaling is the constant for this
geometry, the retention metric and the result for controllers without positions: besides the
bandit (H-A3 supported), URAP gains \numEaGainUrap~dB, and our reading of the heuristic
of~\cite{Mishra2025WCL} reaches \numEaGainWclFast~dB at the shortest delay.

\subsection{E-C per geometry}
\claim{C-EC}Fig.~\ref{fig:F11} and Table~\ref{tab:ablation} give the loss of the optimized RIS
term relative to ideal continuous phase in four settings, per alphabet or defect, with the
large-$N$ closed-form loss. The prior work they agree with concerns discrete and
phase-dependent states~\cite{WuZhang2020Discrete,Abeywickrama2020Practical}; the laboratory
state amplitudes are ASSUMED as in Section~\ref{sec:cases-eb}.

\begin{figure}[t]
\centering
\sremufigure{F11}{\columnwidth}{Loss of the optimized RIS term relative to ideal continuous
phase for each alphabet or defect in the lab-like geometry, the synthetic panel with a near
and a far receiver and the blocked corridor; black ticks: large-$N$ closed form. Data classes:
SIM-sremu (optimized runs), ANALYTIC (ideal alphabets, closed form), ASSUMED (laboratory state
amplitudes), REPORTED (published states of Tang et al.).}{fig:F11}
\end{figure}

% GENERATED by paper/tools/gen_results.py from paper/sweeps/results/E-C/table8.csv.
% Do not edit.
\begin{table}[t]
\caption{Non-ideality ablation (E-C): loss of the optimized RIS term relative to ideal continuous phase, in dB, per geometry, and the large-$N$ closed-form loss $\eta$. SIM-sremu; ideal alphabets ANALYTIC; lab amplitudes ASSUMED; Tang states REPORTED. $|\Gamma_2|$: magnitude of the second laboratory state, the other having unit magnitude; ---: no closed form for stuck cells, whose fixed term depends on the random draw.}
\label{tab:ablation}
\centering
\scriptsize
\setlength{\tabcolsep}{2pt}
\begin{tabular}{@{}lrrrrr@{}}
\toprule
Alphabet or defect & Lab-like & Synth.\ 5\,m & Synth.\ 500\,m & Corr. & $\eta$ \\
\midrule
ideal 1-bit & \ensuremath{-2.89} & \ensuremath{-3.27} & \ensuremath{-3.22} & \ensuremath{-3.28} & \ensuremath{-3.92} \\
ideal 2-bit & \ensuremath{-0.75} & \ensuremath{-0.87} & \ensuremath{-0.89} & \ensuremath{-0.79} & \ensuremath{-0.91} \\
lab states, $|\Gamma_2|$ = 1 & \ensuremath{-2.86} & \ensuremath{-3.18} & \ensuremath{-3.20} & \ensuremath{-3.29} & \ensuremath{-4.10} \\
lab, $|\Gamma_2|$ = 0.7 & \ensuremath{-3.79} & \ensuremath{-4.09} & \ensuremath{-4.17} & \ensuremath{-4.36} & \ensuremath{-5.51} \\
lab, $|\Gamma_2|$ = 0.5 & \ensuremath{-4.46} & \ensuremath{-4.74} & \ensuremath{-4.89} & \ensuremath{-5.15} & \ensuremath{-6.58} \\
Tang RIS1 & \ensuremath{-4.59} & \ensuremath{-4.97} & \ensuremath{-4.91} & \ensuremath{-5.09} & \ensuremath{-5.93} \\
phase-dep.\ amplitude & \ensuremath{-4.14} & \ensuremath{-4.32} & \ensuremath{-4.45} & \ensuremath{-4.66} & \ensuremath{-5.82} \\
lab, 10\,\% dead & \ensuremath{-3.75} & \ensuremath{-4.13} & \ensuremath{-4.09} & \ensuremath{-4.17} & \ensuremath{-5.02} \\
lab, 10\,\% stuck & \ensuremath{-3.60} & \ensuremath{-4.02} & \ensuremath{-3.98} & \ensuremath{-3.98} & --- \\
lab, efficiency 0.5 & \ensuremath{-5.87} & \ensuremath{-6.19} & \ensuremath{-6.21} & \ensuremath{-6.30} & \ensuremath{-7.12} \\
ideal design, lab states & \ensuremath{-3.07} & \ensuremath{-3.63} & \ensuremath{-3.57} & \ensuremath{-3.45} & --- \\
\bottomrule
\end{tabular}
\end{table}


\subsection{FM-band dimensioning (E-D)}\label{sec:cases-ed}
\claim{C-ED}This subsection is model-based dimensioning, not validated. Can a focused RIS
restore an FM broadcast link of \numFmLinkKm~km at \numFmFreqMHz~MHz behind an obstruction?
With a flat-earth two-ray direct path, a coherent knife edge and ground on both RIS hops, the
aperture needed for a RIS path as strong as the direct path is \numFmAreqMin--\numFmAreqMax~m$^2$
over the receiver heights and distances of Fig.~\ref{fig:F12}. At grazing obstruction, over the
same heights and distances, the RIS path of a \numFmPanelMtwo~m$^2$ panel is
\numFmGapMinDb--\numFmGapMaxDb~dB weaker than the direct path. Over all map entries with this direct path, for two placements of the panel and
with or without ground on its hops, the panel suffices in \numFmPrimarySufficient{} of
\numFmPrimaryEntries{}. A four-ray direct path over flat ground, a model-form sensitivity,
gives \numFmFourRaySufficient{} entries. The picture changes with the direct-path model: with the
experimental ITU-R P.526 smooth-earth path of \sremu{}, which is weaker, the same panel suffices in
\numFmPfiveSufficient{} of \numFmPfiveEntries{} entries. That module is not validated, so the
dimensioning question is decided by the direct-path model rather than by the RIS.

\begin{figure}[t]
\centering
\sremufigure{F12}{\columnwidth}{Model-based dimensioning, not validated: aperture that a
focused RIS needs for a path as strong as the direct path, versus the knife-edge loss of the
obstruction, for three RIS-to-receiver distances and two receiver heights; the dotted line
marks the panel size; grey: the experimental P.526 direct path. Data classes: SIM-sremu,
ANALYTIC (panel aperture), ASSUMED (ground constants).}{fig:F12}
\end{figure}

\subsection{E-E: GNU Radio runs}
% GENERATED by paper/tools/gen_results.py from paper/sweeps/results/E-E/summary.json.
% Do not edit.
\begin{table}[t]
\caption{GNU Radio DQPSK modem of stock blocks through the raw request--reply mode (E-E), one fresh run per RIS configuration: focusing, focusing with a seeded fraction $q$ of its cells inverted, uniform (every cell in state 0) and random. AF: isotropic array-factor prediction relative to focusing; gain: least-squares gain of the IQ that GNU Radio received, relative to focusing (AF against gain is an internal consistency check, not a measurement); SNR: estimated from the received IQ over the sample bandwidth; SER: GNU Radio's symbol decisions, 69\,960 scored per configuration; BLER: fraction of the 699 uncoded 200-bit blocks (100 consecutive decisions, no coding, no CRC) with an error. Data classes: SIM-sremu (emulated IQ and the decisions on it), ANALYTIC (AF).}
\label{tab:ee}
\centering
\scriptsize
\setlength{\tabcolsep}{4pt}
\begin{tabular}{@{}lrrrrr@{}}
\toprule
RIS & AF & Gain & SNR & SER & BLER \\
configuration & (dB) & (dB) & (dB) & & \\
\midrule
Focusing & 0.00 & 0.00 & 15.3 & 0 & 0 \\
\quad $q=0.10$ & \ensuremath{-1.98} & \ensuremath{-1.94} & 13.4 & 0 & 0 \\
\quad $q=0.20$ & \ensuremath{-4.43} & \ensuremath{-4.34} & 11.0 & 0 & 0 \\
\quad $q=0.25$ & \ensuremath{-6.20} & \ensuremath{-6.05} & 9.3 & 0 & 0 \\
\quad $q=0.30$ & \ensuremath{-8.01} & \ensuremath{-7.79} & 7.6 & \ensuremath{2.9\times10^{-5}} & \ensuremath{1.4\times10^{-3}} \\
Uniform & \ensuremath{-8.32} & \ensuremath{-8.06} & 7.3 & \ensuremath{5.7\times10^{-5}} & \ensuremath{2.9\times10^{-3}} \\
\quad $q=0.35$ & \ensuremath{-10.95} & \ensuremath{-10.57} & 4.8 & \ensuremath{2.3\times10^{-3}} & 0.11 \\
\quad $q=0.40$ & \ensuremath{-13.83} & \ensuremath{-13.25} & 2.1 & \ensuremath{3.3\times10^{-2}} & 0.78 \\
\quad $q=0.50$ & \ensuremath{-19.02} & \ensuremath{-18.39} & \ensuremath{-3.0} & 0.45 & 1.00 \\
Random & \ensuremath{-23.71} & \ensuremath{-22.73} & \ensuremath{-7.4} & 0.59 & 1.00 \\
\bottomrule
\end{tabular}
\end{table}

\claim{C-EE}The two earlier external-process runs are the frequency-offset sweep between two
processes over the native data plane (V-CFO-3) and the consistency runs of E-A with an external
controller and receiver (VA-1). The GNU Radio DQPSK modem uses \numEeSps{} samples per symbol,
and the partition attenuates the direct path of the E-E scene.
\claim{C-EE-LINK}Table~\ref{tab:ee} lists the static runs of Section~\ref{sec:cases-ee}, one
fresh run with a fresh modem per RIS configuration; the block error rate counts uncoded
200-bit blocks. Where at least one hundred symbol errors were counted with the link up, the
symbol error rate (SER) implies
an implementation loss of \numEeLossMinDb--\numEeLossMaxDb~dB against coherent DQPSK theory at
the estimated $E_s/N_0$; with the link lost, the implied loss is larger.

\claim{C-EE-TIMING}The change point of the switching run is the maximum-likelihood change
point of the received IQ, and in every run every engine and transport counter stayed zero.

\claim{C-EE-SPEED}In \txtEeLaptopSweeps{} sweeps of single runs on a laptop (\txtEeLaptopCpu{}, WSL2,
host load not controlled), the lockstep loop with GNU Radio advanced at an RTF of
\numEeLaptopRtfLoMin--\numEeLaptopRtfLoMax{} at \numEeRateMsps~MS/s and
\numEeLaptopRtfHiMin--\numEeLaptopRtfHiMax{} at \numEeLaptopRateHiMsps~MS/s; on an AWS
\txtEeAwsInstance{} instance it reached \numEeAwsRtfLo{} at \numEeRateMsps~MS/s.

\claim{C-EE-NOTRUN}srsRAN and OpenAirInterface were analysed but not run. srsRAN~4G pads its
transmit stream to the end of the receive block it is about to read and no further, while
lockstep needs the transmit stream one fractional-delay look-ahead beyond the received
samples; a lockstep srsRAN pair therefore needs a per-peer transmit delay and partial receive
replies, and srsRAN~4G has no package for the Ubuntu release used here. The OpenAirInterface
RF simulator would need a new transport for its own TCP framing~\cite{OAIrfsim}. Both are
future work.
